\subsection{相位门与紧李群的绝对定标：秩一体积常数与规范归一化}\label{subsec:unit-circle-phase-compact-lie-normalization}
相位坐标门把 \(\mathbb{U}=\RR/\ZZ\) 识别为相位盒，并在角变量 \(\vartheta=2\pi u\) 下将 Lebesgue 测度 \(\dd u\) 识别为 \(U(1)\) 的归一化 Haar 测度 \(\dd\vartheta/(2\pi)\)（注 \ref{rem:unit-circle-phase-haar-pullback-weight}）。
当该门与有限窗审计接口及 orbit--Levi 刚性并置而不增设额外公理时（见 \cite{Internal2025ResolutionFoldingPhiPi}；并参照 \S\ref{subsubsec:window6-terminal-propositions} 的最小锚点口径），若干通常被视作“规范选择”的连续自由度可被还原为由刚性内生确定的绝对定标问题。
本小节给出一组直接可引用的结构性推论：它们提供由相位门强制的紧李群 Haar 定标，并把离散熵标尺 \(\log(2/\pi)\) 与紧李群 Weyl 推前常数锁定在同一不可调尺度上。
\par

\begin{proposition}[相位门强制的秩一因子尺度刚性：绝对 Haar 体积]\label{prop:phase-gate-rank1-volume-rigidity}
\leanverified{paper\_phase\_gate\_rank1\_volume\_rigidity}
令 \(\vartheta\in[0,2\pi)\) 为相位角，并将其视为 \(U(1)\) 的规范角坐标。
对秩一紧李群因子 \(G\in\{U(1),SU(2),SO(3)\}\)，取一个极大环面嵌入 \(U(1)\hookrightarrow G\)。
称 \(G\) 上的 Riemann 度量为\emph{相位一致的双不变度量}，若该度量双不变且上述 \(U(1)\) 子群上的测地参数与 \(\vartheta\) 一致（从而基本闭轨长度为 \(2\pi\)）。
则该度量在每个 \(G\) 上唯一确定，并且诱导的 Riemann 体积（等价于 Haar 测度的绝对定标）满足
\[
\mathrm{Vol}(U(1))=2\pi,\qquad
\mathrm{Vol}(SU(2))=2\pi^{2},\qquad
\mathrm{Vol}(SO(3))=\pi^{2}.
\]
此外，对射影化相位纤维 \(\mathbb{RP}^{1}=U(1)/\{\pm 1\}\) 有 \(\mathrm{Vol}(\mathbb{RP}^{1})=\pi\)。
\end{proposition}
\begin{proof}
\(U(1)\) 上任意双不变度量必为常曲率圆周度量，且相位一致性要求其基本闭轨长度为 \(2\pi\)，故 \(\mathrm{Vol}(U(1))=2\pi\)。
\par
对 \(SU(2)\)，双不变度量在紧简单情形下仅差一整体缩放；并且在 \(SU(2)\cong S^{3}\) 的识别下，双不变度量对应于半径为 \(r\) 的圆球度量，其体积为 \(\mathrm{Vol}(S^{3}_{r})=2\pi^{2}r^{3}\)。
而极大环面闭轨在该度量下为测地大圆，其长度为 \(2\pi r\)；相位一致性强制该长度为 \(2\pi\)，故 \(r=1\)，从而 \(\mathrm{Vol}(SU(2))=2\pi^{2}\)。
\par
再由 \(SO(3)\cong SU(2)/\{\pm I\}\) 为等距自由二重商，体积减半得 \(\mathrm{Vol}(SO(3))=\pi^{2}\)；同理 \(\mathbb{RP}^{1}=U(1)/\{\pm 1\}\) 为等距二重商，故 \(\mathrm{Vol}(\mathbb{RP}^{1})=\pi\)。证毕。
\end{proof}

\begin{proposition}[相位门选定的 \(u(1)\) 方向消去紧简单因子的尺度自由度]\label{prop:phase-gate-u1-fixes-simple-scale}
\leanverified{paper\_phase\_gate\_u1\_fixes\_simple\_scale}
设 \(\mathfrak g\) 为紧简单李代数，并给定一个由相位门运输得到的一参数子群 \(u\mapsto \exp(uH)\)（其中 \(H\in\mathfrak g\) 在共轭意义下唯一到符号），其在相位角 \(\vartheta=2\pi u\) 下的基本周期为 \(2\pi\)。
则存在唯一的 \(\mathrm{Ad}\)-不变内积 \(\langle\cdot,\cdot\rangle\) 使该 \(U(1)\) 方向满足相位一致归一化（等价地：相应闭轨长度为 \(2\pi\)）。
从而 \(\mathfrak g\) 的双不变度量尺度被完全固定；特别地，与该尺度相容的 Haar 测度绝对定标成为可计算的刚性对象，不含连续可调因子。
\end{proposition}
\begin{proof}
紧简单 \(\mathfrak g\) 上 \(\mathrm{Ad}\)-不变内积空间为一维锥：任取一内积 \(\langle\cdot,\cdot\rangle_{0}\)，任一其它内积必为 \(c\langle\cdot,\cdot\rangle_{0}\)（\(c>0\)）。
相位门指定的生成元 \(H\neq 0\) 确定了一个固定的 \(U(1)\) 方向；相位一致归一化要求该方向的基本闭轨长度为 \(2\pi\)，这等价于对某个固定的 \(H\) 施加一个确定的范数约束 \(\|H\|=\mathrm{const}\)。
由于缩放 \(c\) 仅改变范数的整体尺度，故存在唯一 \(c>0\) 使 \(c\langle\cdot,\cdot\rangle_{0}\) 满足该约束，从而尺度自由度被消去。证毕。
\end{proof}

\begin{corollary}[orbit--Levi 刚性下的 Levi Haar 体积绝对定标]\label{cor:window6-levi-haar-volume-rigidity}
在 window-\(6\) 的 orbit--Levi 口径下，离散维数骨架刚性锁定
\(
\mathfrak g^{\mathrm{Levi}}_6\cong \mathfrak{su}(4)\oplus \mathfrak{su}(2)_L\oplus \mathfrak{su}(2)_R
\)
（推论 \ref{cor:window6-levi-rigidity-pati-salam}），并且 \(\mathfrak{su}(4)\supset \mathfrak{su}(3)\oplus\mathfrak u(1)\) 的相对位置在自同构下共轭唯一（推论 \ref{cor:window6-su4-embedding-rigidity}），且右手块内生选出一维不动点方向 \(\mathfrak u(1)_{\mathrm{eff}}\subset \mathfrak{su}(2)_R\)（推论 \ref{cor:window6-levi-rigidity-sm-residual}）。
将这些由离散证书链锁定的 \(u(1)\) 方向视为相位门可运输对象，并对每个紧简单理想施加命题 \ref{prop:phase-gate-u1-fixes-simple-scale} 的相位一致归一化，
则整个 Levi 群的 Haar 体积形式被离散数据完全确定。
因此，任何以 Haar 体积（或其对数/比值）为输入的维度无关函数，均在该窗口上成为可算法化计算对象，并且不含任何连续拟合参数。
\end{corollary}
\begin{proof}
对 \(\mathfrak{su}(2)_L\) 与 \(\mathfrak{su}(2)_R\)，其 \(\mathrm{Ad}\)-不变内积仅差一缩放，且相位一致闭轨条件固定该缩放；对 \(\mathfrak{su}(4)\) 同理，其缩放由推论 \ref{cor:window6-su4-embedding-rigidity} 所锁定的 \(\mathfrak u(1)\) 方向在命题 \ref{prop:phase-gate-u1-fixes-simple-scale} 下唯一确定。
故 Levi 直和上各理想的绝对定标均被固定，从而 Haar 体积形式亦随之确定。证毕。
\end{proof}

\begin{corollary}[\texorpdfstring{\(\log(2/\pi)\)}{log(2/pi)} 常数与 \(SU(2)\) Weyl 推前归一化的同一性]\label{cor:zigzag-entropy-su2-weyl-constant}
记 \(E_n\) 为 Euler 交错数，并定义离散熵常数
\[
c_{\mathrm{zz}}:=\lim_{n\to\infty}\frac1n\log\frac{E_n}{n!}.
\]
则由定理 \ref{thm:pom-fence-volume-log2pi} 有 \(c_{\mathrm{zz}}=\log(2/\pi)\)。
\par
另一方面，令 \(G=SU(2)\) 取命题 \ref{prop:phase-gate-rank1-volume-rigidity} 的相位一致双不变度量，并以共轭类角参数 \(\phi\in[0,\pi]\) 参数化特征值 \(e^{\pm i\phi}\)。
令 \(\nu\) 为归一化 Haar 概率测度在映射 \(SU(2)\to[0,\pi]\) 下的推前。
则 \(\nu\) 具有密度
\[
\dd\nu(\phi)=\frac{2}{\pi}\sin^{2}\phi\,\dd\phi,
\]
其归一化常数为 \(Z_{SU(2)}=2/\pi\)。
从而得到严格恒等式
\[
c_{\mathrm{zz}}=\log\Bigl(\frac{2}{\pi}\Bigr)=\log Z_{SU(2)}.
\]
\end{corollary}
\begin{proof}
第一部分即定理 \ref{thm:pom-fence-volume-log2pi}。
\par
对第二部分，将 \(SU(2)\cong S^{3}\) 视为单位四元数，取极坐标分解
\(
S^{3}\ni g=(\cos\phi,\ \sin\phi\cdot n)
\)
（其中 \(n\in S^{2}\)），则 \(\phi\in[0,\pi]\) 正是共轭类角参数。
在命题 \ref{prop:phase-gate-rank1-volume-rigidity} 的相位一致定标下，Haar 概率测度对应于 \(S^{3}\) 的均匀体积测度除以 \(\mathrm{Vol}(S^{3})=2\pi^{2}\)。
而 \(S^{3}\) 的体积元在该坐标下分解为
\(
\dd V_{S^{3}}=\sin^{2}\phi\,\dd\phi\,\dd A_{S^{2}}
\)，其中 \(\mathrm{Area}(S^{2})=4\pi\)。
对 \(n\in S^{2}\) 积分并归一化即得
\[
\dd\nu(\phi)=\frac{4\pi}{2\pi^{2}}\sin^{2}\phi\,\dd\phi=\frac{2}{\pi}\sin^{2}\phi\,\dd\phi.
\]
故归一化常数为 \(2/\pi\)，从而结论成立。证毕。
\end{proof}

\begin{remark}[相位一致定标下的体积常数闭包：超越结构仅来自 \texorpdfstring{\(\pi\)}{pi} 与 \texorpdfstring{\(2\)}{2}]\label{rem:phase-gate-volume-constants-closure}
由命题 \ref{prop:phase-gate-rank1-volume-rigidity}，秩一原语集合
\(
\{U(1),SU(2),SO(3),\mathbb{RP}^{1}\}
\)
在相位一致定标下的体积均为 \(q\pi^{k}\) 形（\(q\in\QQ_{>0}\)、\(k\in\{1,2\}\)）。
因此：
\begin{itemize}
  \item 若仅允许有限次直积与有限覆盖/有限商（例如 \(\{\pm 1\}\) 二重商）来生成新的体积常数，则体积仍落在 \(\QQ\cdot\pi^{\ZZ_{\ge 0}}\) 的有限生成子环中；其对数落在 \(\QQ\log\pi+\QQ\log 2\) 的有限维线性张成中。
  \item 若再允许对体积作有限次加法组合（例如体积可加口径下的不交并），则体积仍落在 \(\QQ[\pi]\)；但对数若作用在加性组合上，一般会引入 \(\log(1+\cdots)\) 型非线性项，故需额外约束才能保持在二维张成内。
\end{itemize}
换言之，在本文所采用的相位一致定标与有限代数操作闭包下，出现的超越源点被压缩到 \(\pi\)（及其二重商所引入的 \(2\)）这一对刚性常数。
\end{remark}
